\backmatterentry{参考文献}

\begin{enumerate}
\def\labelenumi{(\Roman{enumi})}
\item
  C. STURM : a) Sur les équations différentielles linéaires du second ordre, Journ. de Math. (1), t. I (1836), p.~106-186 ; b) Sur une classe d'opérations à différences partielles, ibid., p.~373-444.
\item
  J. LIOUVILLE : a) Sur le développement des fonctions ou parties de fonctions en séries dont les divers termes sont assujettis à satisfaire à une même équation différentielle du second ordre contenant un paramètre variable, Journ. de Math. (1), t. I (1836), p.~253-265, t. II (1837), p.~16-35 et 418-436 ; b) D'un théorème dû à M. Sturm et relatif à une classe de fonctions transcendants, ibid., t. I (1836), p.~269-277.
\item
  J. P. GRAM, Ueber die Entwicklung reeller Functionen in Reihen mittelst der Methode der kleinsten Quadrate, J. de Crelle, t. XCIV (1883), p.~41-73.
\item
  H. MINKOWSKI : a) Geometrie der Zahlen, 1 \(^{re}\) éd., Leipzig (Teubner), 1896 ; b) Theorie der konvexen Körper, Gesammelte Abhandlungen, t. II, p.~131-229, Leipzig-Berlin (Teubner), 1911. (Réimpression, New York (Chelsea), 1967.)
\item
  H. POINCARÉ : a) Sur les équations de la Physique mathématique, Rend. Palermo, t. VIII (1894), p.~57-156 (= Œuvres, t. IX, p.~123-196, Paris (Gauthier-Villars), 1954); b) La méthode de Neumann et le problème de Dirichlet, Acta Mathematica, t. XX (1896), p.~59-142 (= Œuvres, t. IX, p.~202-272, Paris (Gauthier-Villars), 1954).
\item
  I. FREDHOLM, Sur une classe d'équations fonctionnelles, Acta Mathematica, t. XXVII (1903), p.~365-390.
\item
  D. HILBERT, Grundzüge einer allgemeinen Theorie der linearen Integralgleichungen, New York (Chelsea), 1953 (= Gött. Nachr., 1904, 1905, 1906, 1910).
\item
  E. SCHMIDT : a) Zur Theorie der linearen und nichtlinearen Integralgleichungen. I. Teil : Entwicklung willkürlicher Funktionen nach Systemen vorgeschriebener, Math. Ann., t. LXIII (1907), p.~433-476 ; b) Ueber die Auflösung linearer Gleichungen mit unendlich vielen Unbekannten, Rend. Palermo, t. XXV (1908), p.~53-77.
\item
  F. RIESZ : a) Untersuchungen über Systeme integrierbarer Funktionen, Math. Ann., t. LXIX (1910), p.~449-497 ; b) Sur certains systèmes singuliers d'équations intégrales, Ann. Ec. Norm. Sup. (3), t. XXVIII (1911), p.~33-62 ; c) Les systèmes d'équations linéaires à une infinité d'inconnues, Paris (Gauthier-Villars), 1913 ; d) Ueber lineare Funktionalgleichungen, Acta Mathematica, t. XLI (1918), p.~71-98 ; e) Zur Theorie des Hilbertschen Raumes, Acta litt. ac scient. (Szeged), t. VII (1934-35), p.~34-38.
\item
  E. HELLY, Ueber Systeme linearer Gleichungen mit unendlich vielen Unbekannten, Monatshefte für Math. und Phys., t. XXXI (1921), p.~60-91.
\item
  H. HAHN, Ueber lineare Gleichungssysteme in linearen Räumen, J. de Crelle, t. CLVII (1927), p.~214-229.
\item
  S. BANACH et H. STEINHAUS, Sur le principe de condensation des singularités, Fund. Math., t. IX (1927), p.~50-61.
\item
  S. BANACH : a) Sur le problème de la mesure, Fund. Math., t. IV (1923), p.~7-33 ; b) Sur les fonctionnelles linéaires, Studia Math., t. I (1929), p.~211-216 et 223-239 ; c) Théorie des opérations linéaires, Warszawa, 1932. (Réimpression, New York (Chelsea), 1963.)
\item
  G. W. MACKEY : a) On infinite-dimensional linear spaces, Trans. Amer. Math. Soc., t. LVII (1945), p.~155-207 ; b) On convex topological spaces, Trans. Amer. Math. Soc., t. LX (1946), p.~519-537.
\item
  G. KÖTHE, Neubegründung der Theorie der vollkommenen Räume, Math. Nachr., t. IV (1951), p.~70-80.
\item
  L. SCHWARTZ, Théorie des distributions, 2 \(^{e}\) édition, Paris (Hermann), 1966.
\item
  J. LINDENSTRAUSS and L. TZAFRIRI, Classical Banach spaces, t. I, Berlin-Heidelberg-New York (Springer), 1977.
\item
  A. GROTHENDIECK : a) Produits tensoriels topologiques et espaces nucléaires, Mem. Amer. Math. Soc., n° 16 (1955); b) Espaces vectoriels topologiques, 3e éd., São Paulo (Publ. Soc. Mat. São Paulo), 1964.
\end{enumerate}

